\section{Discussion}
\label{sec:discussion}
\paragraph{Lessons.}
(1)~Calibrate on programs, not on verdicts: a verdict-level false-discovery rate of a few percent still leaves a wrong clearance in most apps, because the wrong verdicts are spread over apps; per-program calibration bounds what a user of one report cares about.
(2)~State the flow rule in the question: a pricer that reads ``derived data'' narrowly under-prices getters and concentrates wrong clearances on a few widely bundled facts.
(3)~Validity and power separate cleanly: every pricer and every way of combining prices met the target; what a better pricer buys is a larger cleared share.
(4)~A key must be tested against the library code: a learned pricer looks excellent on a key built from the same conventions and fails when the classes it learned from are removed.
(5)~Counting assumptions is not a substitute for pricing them: with no prices, even receipts of one fact are unsafe in 93\% of the apps, and on malware they would clear 22 documented leaks.
(6)~Receipts are worth checking: a handful of facts sits under thousands of verdicts, so a few checks chosen by how many uncleared verdicts they block buy more than a better pricer would, and they also turn priced assumptions into reported leaks.
(7)~A pricer must be judged at the low end: models that answer almost only 0.02 or 0.95 rank well on average but put some real flows at the lowest price, and calibration then refuses them.

\paragraph{What the results mean for practice.}
An analysis designer who ships an unsound taint analysis can ship with it a list of the library assumptions behind each clean verdict and a threshold whose risk was measured on held-out apps, at a price of about \$0.4 to \$16 per corpus, and a short list of facts whose checking would clear the most.
The honest limit is that the certificate speaks about library holes, not about the rest of the analysis (Section~\ref{sec:rq5}).

\section{Threats to Validity}
\label{sec:threats}
\emph{Answer key.}
The main key contains language-model labels; K2 is code-derived but partial, uses FlowDroid's summaries for JDK leaf methods, and is not sound; the audit is AI-assisted except for 67 facts read by one person, and it checks the facts behind the clearances, not the rejected verdicts.
\emph{Exchangeability.}
Apps that bundle the same library are not independent; Section~\ref{sec:rq7} reports grouped splits.
\emph{Corpus.}
Small F-Droid apps, with half of the large ones failing to build, and 33 TaintBench malware apps; no Play Store apps. The transfer to malware rests on 31 apps, too few to estimate a risk near 5\% precisely; it shows no sign of failure, not a second guarantee.
\emph{The analysis.}
Section~\ref{sec:rq5} shows it differs from FlowDroid in both directions; the guarantee is about library holes relative to our analysis.
\emph{Pricers.}
Every pricer is a language model or a classifier trained on FlowDroid's labels; the main pricer is proprietary, and the panel (Section~\ref{sec:rq3}) includes open-weight models. Language-model answers obtained through hosted services may change over time; all prices are stored in the artifact. A few methods were priced through a different access path to the same model (152 of 17{,}830 for one panel model).
\emph{Checks.}
In RQ9 the key plays the person who checks facts, so the gain assumes checks that agree with the key; a person disagrees with it on some facts (Section~\ref{sec:rq4}).
\emph{Definition.}
Results depend on the flow definition; the one choice that mattered (reference keeping) is reported under both readings.
\emph{Construct.}
A ``wrong clearance'' is relative to the key; no key is ground truth.
